% Ad Familiares 11.17 --- headnote
% ad-familiares-11-17, 15 July 43 BC, Rome

\textit{Cicero to D. Brutus, consul-designate, from Rome around
15 July 43 BC. The Perseus dateline gives \textit{Scr. ibidem
eodem tempore quo ep. 16} --- written at the same place and the
same time as 11.16 --- so it shares that letter's dating, set
against the bare year placeholder of the running order. The two
form a pair: both press the candidacy of L. Aelius Lamia, a
leading \textit{eques} whose claims on Cicero went back to 58
BC, and both ask Brutus to throw the weight of the equestrian
centuries, which he controlled, behind him in the praetorian
elections.}

\textit{Where 11.16 opens with its delicate parable about the
right moment to deliver a letter and then makes the bolder
claim --- ``I am the one standing for the praetorship'' ---
this companion note is plainer and more compact: a brief
testimonial to Lamia's intimacy with Cicero and to his public
standing, the worry that the canvassing has grown unruly, and
the direct request that Brutus put all his resources and zeal
into the cause. The repetition (``with all your resources, with
all your zeal'') and the closing entreaty give the short letter
its urgency. With 11.16, this is among the last of Cicero's
letters to D. Brutus to survive; both men would be dead before
the year was out.}
